\documentclass[11pt]{article}

% ---- Encodage et langue ----
\usepackage[T1]{fontenc}
\usepackage[utf8]{inputenc}
\usepackage{lmodern}
\usepackage[french]{babel}

% ---- Mise en page ----
\usepackage[a4paper,margin=2.5cm]{geometry}
\usepackage{setspace}
\onehalfspacing
\usepackage{microtype}
\usepackage{parskip}

% ---- Mathematiques ----
\usepackage{amsmath,amssymb}

% ---- Tableaux (house style) ----
\usepackage{tabularray}
\UseTblrLibrary{booktabs}
\usepackage{booktabs}
\usepackage{pdflscape}
\usepackage{float}
\usepackage{caption}
\captionsetup{labelsep=period}

% ---- Figures natives pgfplots ----
\usepackage{pgfplots}
\pgfplotsset{compat=1.18}
\usepgfplotslibrary{groupplots}
\usepackage{graphicx}

% ---- References croisees ----
\usepackage{hyperref}
\hypersetup{colorlinks=true,linkcolor=black,citecolor=black,urlcolor=blue!50!black}
\usepackage{cleveref}
\usepackage[round]{natbib}
\bibliographystyle{plainnat}

% ===================================================================
%  PALETTE GRAPHIQUE : sens fixe d'une couleur, tenu dans tout le papier.
%    franceblue   : la France, le pays traite, partout.
%    controlgrey  : les pays de comparaison et les placebos, partout.
%    hausse2013   : le choc de hausse d'impot de 2013.
%    baisse2018   : le choc de baisse d'impot de 2018, le prelevement forfaitaire unique.
%  Les deux chocs gardent leur couleur sur toutes les figures ou ils apparaissent.
% ===================================================================
\usepackage{xcolor}
\definecolor{franceblue}{RGB}{31,78,120}
\definecolor{controlgrey}{RGB}{140,140,140}
\definecolor{hausse2013}{RGB}{192,86,33}
\definecolor{baisse2018}{RGB}{33,122,101}
\definecolor{zeroline}{RGB}{110,110,110}

% Style d'axe commun a toutes les figures.
\pgfplotsset{
  paperaxis/.style={
    width=12cm, height=7.2cm,
    axis lines=left,
    tick align=outside,
    tick style={black},
    label style={font=\small},
    tick label style={font=\footnotesize},
    legend style={font=\footnotesize, draw=none, fill=none},
    every axis plot/.append style={line width=0.9pt},
    grid=none,
  },
}

% ---- Note de tableau et de figure sous la legende ----
\newcommand{\tablenotes}[1]{\par\vspace{3pt}{\footnotesize\textit{Notes.} #1}}
\newcommand{\fignote}[1]{\par\vspace{2pt}\begin{minipage}{0.92\linewidth}{\footnotesize\textit{Notes.} #1}\end{minipage}}

% ---- Separateur de milliers a la francaise dans le texte ----
\frenchspacing

\title{\textbf{Après la flat tax : le capital sort plus vite qu'il ne rentre}}
\author{Geoffrey Ducournau\thanks{Contact : \texttt{G.ducournau.voisin@gmail.com}.}}
\date{Octobre 2026}

\begin{document}
\maketitle

\begin{abstract}
\noindent
La fiscalité des revenus du capital est censée orienter l'épargne vers le financement
productif, mais on observe mal où va l'argent qu'une baisse d'impôt libère. Le prélèvement
forfaitaire unique instauré en France en 2018, une imposition proportionnelle des dividendes et
des plus-values à 30 \%, offre un cas net. Nous comparons les entreprises françaises à
celles de cinq économies comparables, avant et après la réforme, sur un panel de plus de
16 millions d'années-sociétés tiré d'Orbis, complété par les sociétés cotées de Compustat
Global et par les comptes nationaux d'Eurostat. La réforme n'a pas fait entrer de capital neuf
dans les fonds propres des sociétés, ni privées ni cotées, et n'a provoqué ni substitution par la
dette ni hausse de l'investissement. L'argent est sorti vers les ménages sous forme de dividendes,
en hausse de 32 \% l'année de la réforme, puis s'est logé en épargne liquide et en
portefeuille d'actions, pas dans le capital des entreprises ni dans l'immobilier. En plaçant côte à
côte la hausse d'impôt de 2013 et la baisse de 2018 dans un cadre unique, nous mettons au jour une
asymétrie. La hausse de 2013 avait réalloué du capital vers les bénéfices retenus ; la baisse
symétrique de 2018 n'en ramène qu'une fraction, moins de la moitié une fois neutralisée
l'anticipation. Le capital part plus facilement qu'il ne revient. Ce cliquet fiscal, cohérent sur
la distribution et sur l'allocation du capital, montre qu'une baisse d'impôt ne défait pas les
déformations d'une hausse passée, un résultat qui importe pour toute évaluation de la fiscalité du
capital.
\end{abstract}


\clearpage
\section{Introduction}

La question de savoir si une baisse de la fiscalité du capital finance l'économie réelle ou
alimente simplement la distribution aux actionnaires structure l'évaluation des réformes fiscales
depuis deux décennies. La théorie des nouvelles vues de l'impôt sur les dividendes prédit que la
fiscalité de la distribution n'affecte pas l'investissement financé par les bénéfices retenus,
tandis que les anciennes vues prédisent un effet sur le coût du capital et donc sur
l'investissement \citep{auerbach1979,poterba1985,chetty2005}. L'arbitrage empirique entre ces
positions reste ouvert, et il dépend de ce que les entreprises et les ménages font effectivement de
l'argent qu'une réforme déplace \citep{hanlon2010}.

Le prélèvement forfaitaire unique, appelé flat tax dans le débat public, a été instauré par
l'article 28 de la loi de finances pour 2018, portée par le gouvernement issu de l'élection
présidentielle de 2017 \citep{loidefinances2018}. Il a remplacé l'imposition au barème progressif
de l'impôt sur le revenu par un taux proportionnel de 30 \%, soit 12,8 \% au titre de l'impôt sur le
revenu et 17,2 \% au titre des prélèvements sociaux, appliqué aux dividendes, aux intérêts et aux
plus-values de cession des personnes physiques résidentes. L'objectif affiché par le gouvernement
était double : simplifier une fiscalité jugée illisible et abaisser son taux marginal pour attirer
l'épargne vers le financement des entreprises et améliorer l'attractivité du pays. Le comité
d'évaluation des réformes de la fiscalité du capital, mandaté par les services du Premier ministre,
a suivi cette réforme et documenté la forte remontée des dividendes versés qui l'a accompagnée, sans
trancher sur l'effet réel \citep{francestrategie2023}. Ce choix d'une imposition proportionnelle du
capital n'est pas isolé en Europe. Il reprend la logique de l'impôt dual des pays nordiques, qui
taxent le revenu du capital à un taux proportionnel et le revenu du travail au barème progressif
depuis les réformes des années quatre-vingt-dix \citep{sorensen1994}, logique que l'Allemagne a
adoptée en 2009 avec son prélèvement libératoire, l'\textit{Abgeltungsteuer}, sur les revenus de
capitaux mobiliers. La réforme française met donc à l'épreuve, sur un cas net et récent, une
doctrine fiscale largement diffusée.

Ce papier pose la question de recherche suivante :

\begin{quote}
\itshape
Une baisse de la fiscalité des dividendes fait-elle entrer du capital neuf dans les entreprises, ou
se contente-t-elle de faire sortir des dividendes vers les ménages ?
\end{quote}

\noindent À cette question principale s'ajoute une question seconde, que la littérature a peu
traitée parce qu'elle exige de comparer deux chocs de sens opposé dans un même cadre : la réponse de
la distribution et de l'allocation du capital à l'impôt est-elle symétrique, de sorte qu'une baisse
défasse l'effet d'une hausse d'ampleur comparable, ou bien existe-t-il un cliquet, le capital
partant plus facilement qu'il ne revient ?

La France offre le terrain qui permet de répondre à ces deux questions, parce qu'elle a parcouru la
même marge fiscale dans les deux sens à cinq ans d'intervalle. Le prélèvement forfaitaire unique de
2018 a abaissé l'imposition des dividendes des personnes physiques résidentes, alors qu'en 2013 la
France avait suivi le chemin inverse en soumettant ces mêmes dividendes au barème progressif et en
alourdissant leur imposition. Ces deux réformes de sens opposé sur la même marge fournissent le
contraste dont la mesure d'un cliquet a besoin, là où les études antérieures n'observent qu'un choc
isolé.

La littérature nous arme pour lire ces chocs et pour situer notre apport. Sur le choc de 2013,
\citet{bach2023} montrent une contraction marquée des dividendes versés par les sociétés françaises,
concentrée sur les sociétés contrôlées par un actionnaire personne physique, et un déplacement de
l'arbitrage vers la rétention des bénéfices plutôt que vers l'investissement. Une première étude du
même choc avait rapporté une hausse de l'investissement \citep{boissel2022}, mais elle a été
rétractée après la mise en évidence de tendances pré-existantes entre sociétés traitées et de
contrôle \citep{bachcomment2023} ; l'ancrage robuste porte donc sur la rétention des bénéfices, non
sur l'investissement. Sur les baisses d'impôt sur les dividendes, les travaux de \citet{chetty2005}
sur la réforme américaine de 2003 documentent une hausse immédiate des distributions, tandis que
\citet{yagan2015} montre que la même baisse n'a relevé ni l'investissement ni l'emploi, un écart
entre la réponse de la distribution et celle de l'économie réelle qui annonce notre propre résultat.
\citet{alstadsater2017} établissent plus largement que la fiscalité de la distribution déplace
l'allocation du capital entre firmes sans nécessairement en changer le niveau agrégé, et
\citet{becker2013} montrent sur un panel international que les impôts sur les versements réorientent
l'investissement vers les firmes qui autofinancent leur croissance. Cette littérature converge vers
un même diagnostic : la fiscalité des dividendes agit d'abord sur le calendrier et la destination
des versements, non sur le niveau du capital productif. Notre contribution s'en distingue sur un
point. Là où chacun de ces travaux étudie un choc isolé, nous posons les deux chocs français dans un
cadre d'estimation unique et nous lisons l'écart de réponse, ce qui fait de l'asymétrie elle-même
l'objet de mesure, et non un sous-produit.

Nous apportons trois résultats. Premièrement, la baisse d'impôt de 2018 ne fait entrer aucun
capital neuf dans les fonds propres des sociétés, ni privées ni cotées, et ne soutient ni
l'investissement ni un recours accru à la dette. Deuxièmement, l'argent libéré sort vers les
ménages, dont les dividendes reçus bondissent de 32 \% l'année de la réforme, et se loge en épargne
liquide et en actions de portefeuille plutôt qu'en financement d'entreprise ou en immobilier.
Troisièmement, la comparaison des deux chocs révèle un cliquet, la baisse de 2018 ne défaisant qu'une
fraction de la réallocation de capital opérée par la hausse de 2013. Nous appuyons ces résultats par
un contrôle positif, la réplication de la hausse de la rétention des bénéfices postérieure à 2013
documentée par \citet{bach2023}, et par un contrôle négatif, l'absence de réponse là où le canal ne
peut pas jouer, les sociétés cotées à l'actionnariat dispersé.

La suite est organisée ainsi. La section~\ref{sec:cadre} développe le cadre théorique et les
hypothèses à partir d'un modèle d'arbitrage de l'actionnaire. La section~\ref{sec:donnees} décrit
les données et la construction des variables. La section~\ref{sec:modele} présente le modèle
empirique. La section~\ref{sec:resultats} rapporte les résultats, hypothèse par hypothèse. La
section~\ref{sec:robustesse} rassemble les tests de robustesse. La section~\ref{sec:conclusion}
conclut.

\section{Cadre théorique et hypothèses}\label{sec:cadre}

Le mécanisme que nous cherchons à identifier est l'arbitrage, par un actionnaire personne physique
qui contrôle la distribution de sa société, entre se verser un euro de bénéfice immédiatement
imposé et le laisser dans la société pour qu'elle l'investisse. Le cadre théorique est celui des
nouvelles vues de l'impôt sur les dividendes, dans lequel l'investissement financé par les bénéfices
retenus n'est pas affecté par le taux d'imposition des dividendes \citep{auerbach1979,king1977,chetty2005}.
Nous le présentons par un modèle à deux dates qui rend le mécanisme explicite et dont les hypothèses
testables descendent directement.

Un actionnaire détient un euro de bénéfice après impôt sur les sociétés. Il peut le distribuer
aujourd'hui, le dividende étant taxé au taux $\tau_d$, et placer le net au-dehors à un rendement $r$
lui-même imposé au taux $\tau_i$ ; ou bien le laisser dans la société, qui l'investit dans un projet
de rendement marginal $\rho$ et le distribue à la date suivante. La première voie procure, à la date
suivante, une consommation nette de $(1-\tau_d)\bigl(1 + r(1-\tau_i)\bigr)$. La seconde procure
$(1-\tau_d)(1+\rho)$. L'actionnaire est indifférent lorsque les deux voies s'égalisent,
\begin{equation}\label{eq:arbitrage}
  (1-\tau_d)\,(1+\rho) = (1-\tau_d)\,\bigl(1 + r(1-\tau_i)\bigr),
\end{equation}
où $\tau_d$ est le taux d'imposition des dividendes, $\tau_i$ celui des autres revenus du capital de
l'actionnaire, $r$ le rendement de son placement extérieur et $\rho$ le rendement marginal du projet
financé par la société. Le taux $\tau_d$ apparaît en facteur des deux membres et s'élimine, de sorte
que le rendement marginal qui déclenche l'investissement vaut
\begin{equation}\label{eq:neutralite}
  \rho^{\star} = r\,(1-\tau_i),
\end{equation}
indépendamment de $\tau_d$. C'est le résultat de neutralité des nouvelles vues : un euro de bénéfice
retenu supporte la taxe sur les dividendes qu'il soit distribué aujourd'hui ou demain, la taxe se
capitalise des deux côtés et ne déplace pas la marge d'investissement. Une variation de $\tau_d$ ne
change donc ni l'entrée de capital neuf, qui ne serait utile que pour financer un investissement
que $\rho^{\star}$ ne commande pas, ni le besoin de financement externe par la dette. Elle ne joue
que sur le calendrier et le niveau des versements. De là découlent les deux premières hypothèses.

\medskip
\noindent\textit{Hypothèse H1. Une baisse de la fiscalité des dividendes ne fait pas entrer de
capital neuf dans les fonds propres des sociétés.}
\smallskip

\noindent Le mécanisme est celui de l'équation~\eqref{eq:neutralite} : la réforme abaisse le coût de
sortie des fonds sans toucher le rendement $\rho^{\star}$ qui commande l'investissement, donc sans
créer de besoin de capital neuf. Le signe attendu sur l'entrée de capital est nul.

\medskip
\noindent\textit{Hypothèse H2. Une baisse de la fiscalité des dividendes ne provoque ni substitution
par la dette ni baisse de l'investissement.}
\smallskip

\noindent Le mécanisme est que, si rien ne quitte le financement interne pour un investissement
supplémentaire, les sociétés n'ont pas à remplacer les fonds distribués par de la dette. Les signes
attendus sur le levier et sur l'investissement sont nuls.

La troisième hypothèse est l'apport central et exige d'enrichir le modèle statique d'une dynamique
d'ajustement. \citet{bach2023} établissent que la hausse d'impôt de 2013 a accru la rétention des
bénéfices des sociétés contrôlées par leur actionnaire, un déplacement vers l'épargne interne plutôt
que vers l'investissement. Notons $R_t$ le stock de capital retenu dans la société à la date $t$ et
$R^{\star}(\tau_d)$ sa cible, croissante en $\tau_d$ puisqu'un impôt plus lourd sur la sortie incite
à conserver les fonds dans la société. Le stock s'ajuste vers sa cible selon
\begin{equation}\label{eq:ajustement}
  R_t - R_{t-1} = \lambda^{\pm}\,\bigl(R^{\star}(\tau_d) - R_{t-1}\bigr),
  \qquad \lambda^{+} \ge \lambda^{-},
\end{equation}
où $R_t$ est le capital retenu, $R^{\star}(\tau_d)$ sa cible fonction du taux d'imposition des
dividendes, et $\lambda^{\pm}$ la vitesse d'ajustement, qui prend la valeur $\lambda^{+}$ quand la
cible monte, après une hausse d'impôt qui pousse à retenir, et $\lambda^{-}$ quand elle baisse, après
une baisse d'impôt qui permet de relâcher. Les coûts d'ajustement et un motif de précaution rendent
la rétention plus rapide que la libération, soit $\lambda^{+} \ge \lambda^{-}$, ce qui engendre un
cliquet : à chocs d'ampleur égale et de sens opposé, le capital retenu réagit davantage à la hausse
qu'à la baisse. La même asymétrie se transmet à la distribution, miroir de la rétention.

\medskip
\noindent\textit{Hypothèse H3. La réponse de la distribution et de l'allocation du capital à la
fiscalité des dividendes est asymétrique : la baisse de 2018 ne défait qu'une fraction de l'effet de
la hausse de 2013, de sorte que le capital part plus facilement qu'il ne revient.}
\smallskip

\noindent Le mécanisme est celui de l'équation~\eqref{eq:ajustement} avec $\lambda^{+} > \lambda^{-}$.
Le signe attendu est une semi-élasticité de la distribution à la baisse inférieure, en valeur
absolue, à celle de la hausse, et une réallocation du capital plus faible au retour qu'à l'aller.

Ce même cadre oriente l'identification. L'arbitrage de l'équation~\eqref{eq:arbitrage} suppose un
actionnaire résident qui pilote la distribution, ce qui ne vaut que dans les sociétés privées
contrôlées. Dans une société cotée à l'actionnariat dispersé et international, aucun contribuable
résident ne fixe la politique de versement, qui obéit au signal de marché et non au taux marginal
$\tau_d$ d'un ménage. La réponse à la réforme ne peut donc vivre que dans les sociétés privées, ce
qui fournit le contrôle négatif des sociétés cotées.

\section{Données et construction des variables}\label{sec:donnees}

Les sociétés privées proviennent d'Orbis, base de données du vendeur Bureau van Dijk, consultée via
Wharton Research Data Services, pour les comptes annuels de 2013 à 2022, étendus à 2010 pour
l'analyse des deux chocs. Les sociétés cotées proviennent de Compustat Global, du vendeur
S\&P Global Market Intelligence, même période et même accès. Le cadrage au niveau des ménages
provient des comptes de secteur et des comptes financiers d'Eurostat, vendeur Office statistique de
l'Union européenne, en accès ouvert.

L'échantillon de sociétés privées retient la France et cinq contrôles comparables : l'Allemagne,
l'Italie, l'Espagne, la Belgique et les Pays-Bas. Une société entre dans le panel équilibré si elle
déclare un capital social au moins 6 années sur la fenêtre. Une observation par société et par année
est conservée, le dépôt au plus grand total de bilan en cas de doublon. Le panel principal compte
16\,177\,643 années-sociétés portées par 2\,138\,541 sociétés. Les variables continues sont
winsorisées aux 1\textsuperscript{er} et 99\textsuperscript{e} centiles. L'échantillon coté compte
24\,686 années-sociétés sur les sociétés cotées des six pays.

La variable expliquée de l'hypothèse H1 est la \textit{variation du capital}, notée $g_{it}$ pour la
société $i$ l'année $t$, variation annuelle du logarithme du capital souscrit et libéré $K_{it}$,
\begin{equation}\label{eq:var-capital}
  g_{it} = \ln(1 + K_{it}) - \ln(1 + K_{i,t-1}).
\end{equation}
Nous lui adjoignons la \textit{hausse du capital}, indicatrice $a_{it}$ valant un si le capital
souscrit a augmenté d'une année sur l'autre et zéro sinon. Pour l'hypothèse H2, le \textit{levier}
$\ell_{it}$ est le total des dettes $L_{it}$ rapporté au total de l'actif $A_{it}$,
\begin{equation}\label{eq:levier}
  \ell_{it} = L_{it} / A_{it},
\end{equation}
la \textit{croissance de la dette} est la variation du logarithme de la dette financière, dette à
long terme plus emprunts courants, et l'\textit{investissement} est la variation du logarithme des
immobilisations corporelles. Pour la source des fonds et pour l'hypothèse H3, les \textit{réserves
sur actif} $r_{it}$ sont les capitaux propres moins le capital souscrit, rapportés au total de
l'actif,
\begin{equation}\label{eq:reserves}
  r_{it} = (E_{it} - K_{it}) / A_{it},
\end{equation}
où $E_{it}$ est le total des capitaux propres, et le \textit{ratio de fonds propres} $e_{it}$ est
\begin{equation}\label{eq:fonds-propres}
  e_{it} = E_{it} / A_{it}.
\end{equation}
Au niveau coté, l'\textit{émission de capital} est la trésorerie levée par vente d'actions rapportée
à l'actif. Au niveau des ménages, la \textit{distribution} $Y_{ct}$ est le logarithme des dividendes
distribués reçus par les ménages du pays $c$ l'année $t$.

Le modèle ne retient qu'une variable de contrôle au niveau de la société, la \textit{taille}
$s_{it}$, logarithme du total de l'actif,
\begin{equation}\label{eq:taille}
  s_{it} = \ln(1 + A_{it}).
\end{equation}
Les effets fixes sont de deux ordres : un effet fixe de société $\alpha_i$, qui absorbe toute
différence de niveau stable entre sociétés, et un effet fixe d'année $\gamma_t$, qui absorbe tout
choc annuel commun aux six pays. L'âge de la société n'est pas introduit comme contrôle : étant la
différence exacte entre l'année courante et l'année de constitution, il est déjà absorbé par la
combinaison des effets fixes de société et d'année. Les erreurs sont regroupées au niveau du pays,
l'unité à laquelle le traitement est affecté. Les définitions complètes, les champs source et les
formules figurent dans la table~\ref{table:definitions-variables} en annexe, que la
table~\ref{table:stats-descriptives} des statistiques descriptives accompagne.

\section{Stratégie d'identification et modèle empirique}\label{sec:modele}

La stratégie d'identification repose sur un fait institutionnel : la France a parcouru la même marge
fiscale, l'imposition des dividendes des personnes physiques résidentes, dans les deux sens à cinq
ans d'intervalle. La loi de finances 2013 a ramené les dividendes au barème progressif et alourdi
leur imposition, une hausse que nous estimons sur une fenêtre avant de 2010 à 2012 et une fenêtre
après de 2013 à 2016. La loi de finances 2018, en son article 28, a instauré le prélèvement
forfaitaire unique à 30 \%, une baisse que nous estimons sur une fenêtre avant de 2014 à 2017 et une
fenêtre après de 2018 à 2022. La France est le seul pays à offrir ce couple de chocs symétriques et
rapprochés, ce qui fournit le contraste dont la mesure du cliquet a besoin. Les cinq pays de
contrôle, l'Allemagne, les Pays-Bas, l'Espagne, la Belgique et l'Italie, servent de contre-factuel
sur toute la fenêtre 2010 à 2022 : chacun dispose d'un régime forfaitaire du capital instauré avant
la fenêtre, l'\textit{Abgeltungsteuer} allemande à 26,375 \% en 2009, le \textit{Box 3} néerlandais
en 2001, la base de l'épargne espagnole en 2007, le précompte mobilier belge et la \textit{cedolare
secca} italienne, et aucun ne reproduit l'aller-retour hausse puis baisse de la France. Deux d'entre
eux ont néanmoins ajusté leur taux forfaitaire dans la fenêtre, l'Italie portant la \textit{cedolare
secca} à 26 \% en 2014 et la Belgique son précompte mobilier à 30 \% en 2017, sans parcourir la
marge dans les deux sens ; la permutation placebo, qui classe la France parmi les six pays, et la
robustesse au choix du pool de comparaison traitent cette composition. Le
tableau~\ref{table:identification}, en annexe, récapitule cette variation identifiante. Les dates
des textes de loi sont à confirmer avant diffusion.

L'estimateur principal est une double différence à effets fixes bidirectionnels, dite TWFE pour
\textit{two-way fixed effects}. La France est le seul groupe traité et les cinq pays comparables forment le
groupe de contrôle ; chaque réforme a une date de mise en \oe uvre unique. La spécification compare
donc un seul groupe traité à un ensemble de contrôles autour d'une seule date de traitement, ce qui
en fait un schéma à deux périodes et deux groupes généralisé à plusieurs années. Pour une variable
expliquée $y_{it}$,
\begin{equation}\label{eq:did}
  y_{it} = \beta \, (D_c \cdot P_t) + \theta \, s_{it} + \alpha_i + \gamma_t + \varepsilon_{it}.
\end{equation}
Dans cette équation, $y_{it}$ est la variable expliquée de la société $i$ l'année $t$ ; $D_c$ vaut
un si le pays $c$ est la France, le pays traité, et zéro pour un pays de contrôle ; $P_t$ vaut un à
partir de l'année de la réforme et zéro avant ; $s_{it}$ est le logarithme du total de l'actif ;
$\alpha_i$ est l'effet fixe de société ; $\gamma_t$ l'effet fixe d'année ; et $\varepsilon_{it}$ le
résidu. Le coefficient d'intérêt $\beta$ mesure l'effet sur la France après la réforme, au-delà de
la trajectoire commune des contrôles ; son signe attendu suit les hypothèses, nul pour l'entrée de
capital et les marges de substitution, positif puis plus faible pour les marges de rétention au
retour qu'à l'aller.

Le choix du TWFE appelle une justification, car cet estimateur a fait l'objet de critiques récentes.
Lorsque le traitement est échelonné dans le temps sur plusieurs groupes, la double différence à
effets fixes mélange des comparaisons entre unités traitées tôt et tard et peut pondérer certains
effets négativement \citep{dechaisemartin2020,goodmanbacon2021}. Ces critiques ne s'appliquent pas
ici : un seul groupe est traité et chaque choc a une date unique, de sorte qu'aucune comparaison
interdite entre cohortes de traitement n'intervient et que $\beta$ s'interprète sans ambiguïté comme
la double différence entre la France et ses contrôles avant et après la réforme.

La dynamique est lue par une étude d'événement qui remplace l'interaction unique par des
interactions année par année, suivant la spécification à période de référence omise de
\citet{sunabraham2021},
\begin{equation}\label{eq:event}
  y_{it} = \sum_{j \neq j_0} \tau_j \, \bigl(D_c \cdot \mathbf{1}\{t = j\}\bigr)
           + \theta \, s_{it} + \alpha_i + \gamma_t + \varepsilon_{it},
\end{equation}
où $\tau_j$ est l'écart de la France au pays de contrôle l'année $j$, $\mathbf{1}\{\cdot\}$
l'indicatrice, et $j_0$ l'année de référence omise, 2017 pour le choc de 2018 et 2012 pour le choc
de 2013. Chaque coefficient $\tau_j$ se lit comme l'écart France contre contrôles l'année $j$
relativement à l'année de référence. La platitude des $\tau_j$ avant la réforme teste l'hypothèse de
tendances parallèles ; une rupture après la réforme en signe l'effet dynamique.

L'inférence ne repose pas sur l'erreur-type regroupée, car le nombre de pays traités est de un et le
nombre de groupes de six, trop faible pour une inférence asymptotique fiable. Nous reportons la
statistique de Student à titre indicatif, et l'inférence s'appuie sur une permutation placebo de
l'affectation du traitement à chacun des pays, dont la valeur $p$ est le rang de la France parmi les
six pays. Avec six pays, cette valeur $p$ a un plancher mécanique d'un sixième, soit environ
$0{,}167$, ce que nous assumons explicitement. Au niveau agrégé des ménages, l'effet de chaque choc
est également estimé par un contrôle synthétique \citep{abadie2010,abadie2015}, la France étant
reconstruite comme combinaison convexe des donneurs ajustée sur la période antérieure au choc.

\section{Résultats}\label{sec:resultats}

\subsection{La baisse d'impôt ne fait entrer aucun capital neuf}

Le tableau~\ref{table:h1-capital} rapporte l'effet de la réforme de 2018 sur trois marges de
l'entrée de capital : la variation du capital souscrit en colonne~(1), la probabilité d'une hausse
du capital en colonne~(2) et le ratio de fonds propres en colonne~(3). Sur la première marge, le
coefficient de double différence vaut $-0{,}0034$, une grandeur économiquement nulle : rapportée au
niveau, elle signifie que le capital souscrit des sociétés françaises croît environ $0{,}34\%$ moins
vite que celui des pays comparables chaque année, un écart négligeable et de signe contraire à l'entrée de
capital qu'une baisse d'impôt attractive était censée provoquer.\footnote{$-0{,}0034$ est un
coefficient en quasi-logarithme, soit un écart de taux de croissance annuel d'environ
$-0{,}34\%$.} La permutation placebo place même la France au plus faible écart des six pays, avec
une valeur $p$ de $1{,}000$, c'est-à-dire que cinq pays non traités affichent un écart plus marqué
que le pays réformé. La probabilité d'une hausse du capital recule elle aussi très légèrement,
$-0{,}0049$, et le ratio de fonds propres, en colonne~(3), bouge de $+0{,}0474$ mais sans aucune
précision, la statistique de Student valant $+0{,}9$ et le placebo plaçant la France au deuxième rang
seulement. L'étude d'événement de la figure~\ref{fig:entree} confirme cette lecture : la trajectoire
de l'écart France contre les pays comparables est plate année après année avant 2018 et ne connaît aucune rupture
ensuite, les barres d'erreur couvrant zéro sur toute la fenêtre. L'hypothèse H1 est soutenue : la
baisse d'impôt ne fait pas entrer de capital neuf dans les fonds propres. Ce résultat n'est pas une
incertitude mais un nul estimé avec précision, cohérent avec la neutralité de
l'équation~\eqref{eq:neutralite} et avec l'absence d'effet réel documentée par \citet{yagan2015}
sur la baisse américaine de 2003.

\begin{figure}[ht]
  \centering
  \begin{tikzpicture}
    \begin{axis}[paperaxis,
      xlabel={Année}, ylabel={Écart France moins comparables (variation du log capital)},
      xtick={2014,2015,2016,2017,2018,2019,2020,2021,2022},
      xticklabel style={/pgf/number format/1000 sep=},
    ]
      \draw[zeroline, line width=0.6pt] (axis cs:2013.5,0) -- (axis cs:2022.5,0);
      % La France, en bleu, avec intervalle a 95 pour cent.
      \addplot[franceblue, mark=*, mark options={scale=0.8}, error bars/.cd,
        y dir=both, y explicit, error bar style={franceblue}]
        table[col sep=comma, x=year, y=coef, y error expr=1.96*\thisrow{std_error}]
        {figures/entree_capital.csv};
      % Reforme de 2018 : marqueur vert, la couleur de la baisse d'impot.
      \draw[baisse2018, dashed, line width=1pt]
        ({axis cs:2017.5,0}|-{rel axis cs:0,0}) -- ({axis cs:2017.5,0}|-{rel axis cs:0,1});
    \end{axis}
  \end{tikzpicture}
  \caption{Étude d'événement de l'entrée de capital autour de la réforme de 2018}
  \label{fig:entree}
  \fignote{Coefficients $\tau_j$ de l'équation~\eqref{eq:event} pour la variation du logarithme du
  capital souscrit, écart de la France aux cinq pays comparables l'année $j$, année de référence 2017 omise.
  Barres d'erreur à quatre-vingt-quinze pour cent sur l'erreur-type regroupée par pays, indicative.
  La ligne verte verticale marque la réforme de 2018. Un profil plat avant et après la réforme
  indique l'absence d'entrée de capital neuf.}
\end{figure}


{
\begin{table}[ht]
    \centering
    \caption{Entrée de capital après la réforme de 2018}\label{table:h1-capital}
    \tablenotes{Double différence de l'équation~\eqref{eq:did}, France contre les cinq pays comparables,
    effets fixes de société et d'année, contrôle de taille, période 2013 à 2022. La ligne
    \textit{France $\times$ après 2018} donne le coefficient $\beta$ ; la statistique de Student
    regroupée par pays est entre parenthèses. Les étoiles marquent la significativité d'après la
    valeur $p$ bilatérale associée à cette statistique : *** pour $p < 0{,}001$, ** pour
    $p < 0{,}01$, * pour $p < 0{,}05$. Le nombre de groupes étant de six, ces étoiles et la
    statistique de Student sont indicatives ; l'inférence du papier repose sur l'effet du contrôle synthétique et sur la
    valeur $p$ de permutation placebo, dont le plancher mécanique est $1/6 \approx 0{,}167$ ; une
    valeur de $1{,}000$ signifie que la France a l'écart le plus faible des six pays. Les variables
    continues sont winsorisées aux 1\textsuperscript{er} et 99\textsuperscript{e} centiles. Le ratio
    de fonds propres est calculé sur les années-sociétés à actif strictement positif. Les définitions
    sont dans la table~\ref{table:definitions-variables}.\\}
    \footnotesize
    \begin{tblr}{colspec={l*{3}{X[1,c]}}}
    \toprule
     & (1) & (2) & (3) \\
    Variables expliquées : & Variation du capital & Hausse du capital & Ratio de fonds propres \\
    \midrule
    France $\times$ après 2018 & $-0{,}0034^{***}$ & $-0{,}0049^{***}$ & $+0{,}0474$ \\
                               & $(-9{,}9)$  & $(-5{,}2)$  & $(+0{,}9)$  \\
     & & & \\
    Effet synthétique & $+0{,}0012$ & $-0{,}0059$ & $-0{,}0160$ \\
    Placebo $p$       & $1{,}000$   & $1{,}000$   & $0{,}833$   \\
     & & & \\
    Effets fixes société & Oui & Oui & Oui \\
    Effets fixes année   & Oui & Oui & Oui \\
    Contrôle de taille   & Oui & Oui & Oui \\
    Regroupement         & Pays & Pays & Pays \\
    Observations         & 16\,177\,643 & 16\,177\,643 & 16\,177\,643 \\
    \bottomrule
    \end{tblr}
\end{table}
}

\subsection{Absence de substitution par la dette et d'effet réel}

Le tableau~\ref{table:h2-substitution} rassemble les quatre marges de l'hypothèse de substitution :
le levier en colonne~(1), la croissance de la dette en colonne~(2), l'investissement en colonne~(3)
et la productivité du travail en colonne~(4). Aucune ne montre d'effet dans le sens attendu, et
aucune statistique de Student ne dépasse le seuil usuel. Le résultat le plus informatif tient à la
complémentarité comptable du financement : la part des fonds propres et la part des dettes dans le
bilan somment à l'unité, de sorte que le coefficient sur le levier, $-0{,}0472$, et celui sur le
ratio de fonds propres du tableau précédent, $+0{,}0474$, sont la même grandeur au signe près. S'il
y avait substitution, le point de fonds propres perdu se retrouverait en dette ; or les deux
coefficients sont nuls et imprécis, le placebo du levier plaçant la France au deuxième rang. Pour un
point de part de fonds propres en moins, on ne mesure donc pas un point de dette en plus, mais zéro.
L'investissement, en colonne~(3), affiche $+0{,}0683$ mais avec une statistique de $+1{,}5$ qui ne
franchit aucun seuil, et la productivité ne bouge pas. L'hypothèse H2 n'est pas soutenue au sens où
la substitution prédite par les anciennes vues n'apparaît pas ; conformément à
l'équation~\eqref{eq:neutralite}, les sociétés n'ont pas eu à remplacer des fonds distribués, faute
d'un surcroît d'investissement à financer. Cette absence d'effet réel prolonge le constat de
\citet{alstadsater2017} et de \citet{becker2013}, pour qui la fiscalité des versements déplace
l'allocation du capital sans en relever le niveau productif.

L'étude d'événement de la figure~\ref{fig:levier} confirme ce nul dans le temps. L'écart de levier
entre la France et ses pays comparables, un peu positif en début de période, reflue vers zéro avant
la réforme et ne remonte pas après 2018 ; il devient même légèrement négatif. Ce levier agrège
l'ensemble des dettes au bilan, et son absence de hausse signe l'absence de la substitution de la
dette aux fonds propres prédite par les anciennes vues. La dette financière garantie souscrite
pendant la crise sanitaire fait exception et monte en 2020 à 2022, mais elle est logée en trésorerie
plutôt que convertie en distribution ; ce mouvement est propre au Covid, non à la réforme, et nous
le caractérisons à la section~\ref{sec:robustesse}.

\begin{figure}[ht]
  \centering
  \begin{tikzpicture}
    \begin{axis}[paperaxis,
      xlabel={Année}, ylabel={Écart France moins comparables (levier)},
      xtick={2013,2014,2015,2016,2017,2018,2019,2020,2021,2022},
      xticklabel style={/pgf/number format/1000 sep=, rotate=45, anchor=east},
    ]
      \draw[zeroline, line width=0.6pt] (axis cs:2012.5,0) -- (axis cs:2022.5,0);
      % La France, en bleu, avec intervalle a 95 pour cent.
      \addplot[franceblue, mark=*, mark options={scale=0.8}, error bars/.cd,
        y dir=both, y explicit, error bar style={franceblue}]
        table[col sep=comma, x=year, y=coef, y error expr=1.96*\thisrow{std_error}]
        {figures/levier_event.csv};
      % Reforme de 2018 : marqueur vert, la couleur de la baisse d'impot.
      \draw[baisse2018, dashed, line width=1pt]
        ({axis cs:2017.5,0}|-{rel axis cs:0,0}) -- ({axis cs:2017.5,0}|-{rel axis cs:0,1});
    \end{axis}
  \end{tikzpicture}
  \caption{Étude d'événement du levier autour de la réforme de 2018}
  \label{fig:levier}
  \fignote{Coefficients $\tau_j$ de l'équation~\eqref{eq:event} pour le levier, total des dettes sur
  total de l'actif, écart de la France aux cinq pays comparables l'année $j$, année de référence 2017 omise.
  Barres d'erreur à quatre-vingt-quinze pour cent sur l'erreur-type regroupée par pays, indicative.
  La ligne verte verticale marque la réforme de 2018. Le levier français, un peu supérieur en début
  de période, reflue vers celui des pays comparables et ne remonte pas après la réforme. Ce levier
  agrège l'ensemble des dettes ; la hausse de la dette financière garantie pendant la crise
  sanitaire, logée en trésorerie, n'y apparaît pas et est caractérisée à la
  section~\ref{sec:robustesse}.}
\end{figure}


{
\begin{table}[ht]
    \centering
    \caption{Substitution par la dette et effets réels après 2018}\label{table:h2-substitution}
    \tablenotes{Double différence de l'équation~\eqref{eq:did}, mêmes spécification, période et
    conventions d'inférence que le tableau~\ref{table:h1-capital}. La statistique de Student
    regroupée par pays est entre parenthèses ; les étoiles marquent la significativité d'après la
    valeur $p$ bilatérale, *** pour $p < 0{,}001$, ** pour $p < 0{,}01$, * pour $p < 0{,}05$, et
    sont indicatives vu les six groupes. L'inférence repose sur l'effet synthétique et la valeur $p$
    de permutation placebo. Les définitions sont dans la table~\ref{table:definitions-variables}.\\}
    \footnotesize
    \begin{tblr}{colspec={l*{4}{X[1,c]}}}
    \toprule
     & (1) & (2) & (3) & (4) \\
    Variables expliquées : & Levier & Croissance dette & Investissement & Productivité \\
    \midrule
    France $\times$ après 2018 & $-0{,}0472$ & $+0{,}0302$ & $+0{,}0683$ & $+0{,}0049$ \\
                               & $(-0{,}9)$  & $(+0{,}2)$  & $(+1{,}5)$  & $(+0{,}4)$  \\
     & & & & \\
    Effet synthétique & $+0{,}0165$ & $+0{,}1763$ & $+0{,}0082$ & $-0{,}0896$ \\
    Placebo $p$       & $0{,}833$   & $0{,}333$   & $0{,}667$   & $0{,}500$   \\
     & & & & \\
    Effets fixes société & Oui & Oui & Oui & Oui \\
    Effets fixes année   & Oui & Oui & Oui & Oui \\
    Contrôle de taille   & Oui & Oui & Oui & Oui \\
    Regroupement         & Pays & Pays & Pays & Pays \\
    \bottomrule
    \end{tblr}
\end{table}
}

\subsection{La destination des fonds libérés : les ménages}

L'argent que la réforme ne dirige pas vers le capital des sociétés sort vers les ménages. Les
comptes nationaux montrent que les dividendes reçus par les ménages français augmentent de 32 \%
l'année même de la réforme, passant de 32\,592 à 43\,088 millions d'euros courants entre 2017 et
2018.\footnote{$43\,088 / 32\,592 - 1 = 0{,}322$ ; dividendes distribués reçus par les ménages
français, millions d'euros courants, 2017 et 2018, Eurostat.} La figure~\ref{fig:dividendes} place
la France parmi ses cinq pays comparables en indice base 100 en 2017 et isole nettement la rupture française :
l'indice français saute à 132 en 2018 quand la plus forte hausse parmi les pays comparables, celle des Pays-Bas,
atteint 114, de sorte que la France est le seul pays à rompre franchement l'année de sa réforme.
L'ampleur de ce saut, 32 points d'indice contre au plus 14 chez les pays comparables, mesure directement la
part des fonds qui sort du périmètre des sociétés. Du côté de l'emploi de ces fonds, les comptes
financiers des ménages montrent que le surcroît se loge en dépôts et en achats d'actions cotées sur
le marché secondaire, non en actions non cotées, c'est-à-dire non dans le capital des sociétés
privées, et l'indice des prix du logement français progresse moins que celui des pays comparables, ce qui ne
soutient pas un canal immobilier. L'argent reste donc dans l'épargne liquide et le portefeuille des
ménages ; il ne finance pas l'appareil productif, exactement l'écart entre réponse de la
distribution et réponse réelle que \citet{chetty2005} et \citet{yagan2015} relèvent sur la réforme
américaine de 2003.

Le contraste entre les deux versants du résultat se résume en une phrase. La
figure~\ref{fig:entree} ne montre aucune entrée de capital neuf dans les fonds propres des sociétés
françaises après la réforme de 2018, et le tableau~\ref{table:h2-substitution} aucun surcroît
d'investissement ni de dette ; la figure~\ref{fig:dividendes} montre au contraire une forte sortie
de dividendes vers les ménages, de 32 \% la seule année de la réforme. Autrement dit, après la
baisse d'impôt, rien n'entre dans le capital des entreprises et rien ne s'y investit, tandis qu'une
large part en sort vers les ménages. La baisse d'impôt n'a donc pas financé les entreprises ; elle a
alimenté la distribution.

\subsection{Une réponse asymétrique à la hausse et à la baisse d'impôt}

Nous plaçons enfin la hausse d'impôt de 2013 et la baisse de 2018 dans le même cadre afin de mesurer
l'asymétrie annoncée par l'équation~\eqref{eq:ajustement}. Sur la distribution, mesurée au niveau
agrégé des ménages, la semi-élasticité au choc de 2013 est fortement négative et celle au choc de
2018 positive mais plus petite. Sur une base lissée des trois années antérieures à chaque choc, qui
neutralise le report des versements d'une année isolée, la semi-élasticité vaut $-0{,}40$ pour la
hausse de 2013 et $+0{,}18$ pour la baisse de 2018 : la baisse ne défait donc que 45 \% de l'effet
de la hausse, un peu moins de la moitié.\footnote{$0{,}18 / 0{,}40 = 0{,}45$ ; semi-élasticités sur
base lissée des trois années antérieures, chocs de 2018 et de 2013.} Ce rapport ne tient pas au
choix de la fenêtre de référence : mesurée sur le flux brut d'une année sur l'autre, la
semi-élasticité vaut $-0{,}52$ en 2013 et $+0{,}23$ en 2018, soit le même rapport de 45
\%.\footnote{$0{,}23 / 0{,}52 = 0{,}45$ ; semi-élasticités sur le flux brut de l'année du choc, 2018
et 2013.}

Sur l'allocation du capital au niveau des sociétés privées, le tableau~\ref{table:h3-cliquet} montre
le même sens et chiffre l'asymétrie. La hausse de 2013 a réalloué du capital vers les bénéfices
retenus : la part des réserves dans l'actif monte de $+0{,}0942$ et la part des fonds propres de
$+0{,}1094$, l'une et l'autre nettement, avec des statistiques de Student de $+2{,}2$ et $+2{,}6$. La
baisse de 2018 ne ramène ces mêmes parts que de $+0{,}0332$ et $+0{,}0403$, des grandeurs trois fois
plus faibles et sans précision. La baisse ne défait donc qu'environ un tiers de la réallocation
opérée par la hausse, le capital retenu revenant dans un rapport de $1$ à $2{,}8$ sur les réserves et
de $1$ à $2{,}7$ sur les fonds propres.\footnote{$0{,}0942 / 0{,}0332 = 2{,}84$ pour la rétention,
$0{,}1094 / 0{,}0403 = 2{,}71$ pour les fonds propres ; coefficients de double différence aux chocs
de 2013 et de 2018.} La figure~\ref{fig:cliquet} résume cette asymétrie sur les trois marges, la
distribution, les réserves et les fonds propres, en portant côte à côte l'ampleur de chaque réponse :
la barre de la hausse de 2013 dépasse celle de la baisse de 2018 sur chacune. L'hypothèse H3 est
soutenue sur les deux marges, la distribution et le capital : conformément au cliquet d'ajustement
avec $\lambda^{+} > \lambda^{-}$, le capital part plus facilement qu'il ne revient.

\begin{figure}[ht]
  \centering
  \begin{tikzpicture}
    \begin{axis}[paperaxis,
      xlabel={Année}, ylabel={Dividendes reçus par les ménages (indice, 2017 $=100$)},
      xtick={2013,2014,2015,2016,2017,2018,2019,2020,2021,2022},
      xticklabel style={/pgf/number format/1000 sep=},
      legend pos=north west, legend cell align=left,
      ymin=40,
    ]
      % Cinq pays de comparaison, tous en gris : meme sens de couleur dans tout le papier.
      \addplot[controlgrey] table[col sep=comma, x=year, y=DE] {figures/dividendes_menages.csv};
      \addplot[controlgrey] table[col sep=comma, x=year, y=IT] {figures/dividendes_menages.csv};
      \addplot[controlgrey] table[col sep=comma, x=year, y=ES] {figures/dividendes_menages.csv};
      \addplot[controlgrey] table[col sep=comma, x=year, y=BE] {figures/dividendes_menages.csv};
      \addplot[controlgrey] table[col sep=comma, x=year, y=NL] {figures/dividendes_menages.csv};
      \addlegendentry{Pays de comparaison}
      % La France, toujours en bleu.
      \addplot[franceblue, line width=1.6pt, mark=*, mark options={scale=0.8}]
        table[col sep=comma, x=year, y=FR] {figures/dividendes_menages.csv};
      \addlegendentry{France}
      % Reforme de 2018 : marqueur vert, la couleur de la baisse d'impot.
      \draw[baisse2018, dashed, line width=1pt] (axis cs:2018,40) -- (axis cs:2018,200);
    \end{axis}
  \end{tikzpicture}
  \caption{Dividendes reçus par les ménages, France et cinq pays comparables, indice base 100 en 2017}
  \label{fig:dividendes}
  \fignote{Dividendes distribués reçus par les ménages, poste des comptes de secteur
  (ressources du secteur des ménages), millions d'euros courants, ramenés à un indice valant 100 en
  2017. En bleu la France, le pays traité ; en gris les cinq pays de comparaison, l'Allemagne,
  l'Italie, l'Espagne, la Belgique et les Pays-Bas. La ligne verte verticale marque la réforme de
  2018. Source : Eurostat, Office statistique de l'Union européenne.}
\end{figure}


\begin{figure}[ht]
  \centering
  \begin{tikzpicture}
    \begin{axis}[paperaxis,
      ybar, bar width=14pt,
      ylabel={Réponse à l'impôt (valeur absolue)},
      symbolic x coords={distribution,retention,fondspropres},
      xtick=data,
      xticklabels={Distribution, Réserves, Fonds propres},
      ymin=0,
      enlarge x limits=0.28,
      legend pos=north east, legend cell align=left,
      nodes near coords, nodes near coords style={font=\scriptsize},
      every node near coord/.append style={/pgf/number format/.cd, use comma, fixed, fixed zerofill, precision=2},
    ]
      % La hausse de 2013 en orange, la baisse de 2018 en vert : meme code que partout.
      \addplot[fill=hausse2013, draw=hausse2013!60!black]
        table[col sep=comma, x=marge, y=choc2013] {figures/cliquet.csv};
      \addlegendentry{Choc de hausse, 2013}
      \addplot[fill=baisse2018, draw=baisse2018!60!black]
        table[col sep=comma, x=marge, y=choc2018] {figures/cliquet.csv};
      \addlegendentry{Choc de baisse, 2018}
    \end{axis}
  \end{tikzpicture}
  \caption{Le cliquet fiscal : la réponse à la hausse de 2013 dépasse la réponse à la baisse de 2018}
  \label{fig:cliquet}
  \fignote{Ampleur de la réponse à chaque choc, en valeur absolue, sur trois marges. La distribution
  est la semi-élasticité des dividendes des ménages, sur base lissée des trois années antérieures au
  choc ; la rétention et les fonds propres sont les coefficients de double différence sur les parts
  de l'actif, au niveau des sociétés. Par convention de couleur, l'orange est la hausse d'impôt de
  2013 et le vert la baisse de 2018. Sur les trois marges, la réponse à la hausse de 2013 dépasse la
  réponse à la baisse de 2018 : le capital part plus facilement qu'il ne revient. La semi-élasticité
  de la distribution est négative en 2013 et positive en 2018 ; seules les amplitudes sont tracées.}
\end{figure}


{
\begin{table}[ht]
    \centering
    \caption{Le cliquet du capital : réallocation à la hausse de 2013 et à la baisse de 2018}
    \label{table:h3-cliquet}
    \tablenotes{Double différence de l'équation~\eqref{eq:did} estimée séparément autour de chaque
    choc, France contre les cinq pays comparables, effets fixes de société et d'année, contrôle de taille,
    panel étendu à 2010. La colonne \textit{Choc de 2013} prend l'année 2013 comme début de période
    traitée sur la fenêtre 2010 à 2016 ; la colonne \textit{Choc de 2018} prend 2018 sur la fenêtre
    2015 à 2022. La statistique de Student regroupée par pays est entre parenthèses ; les étoiles
    marquent la significativité d'après la valeur $p$ bilatérale, *** pour $p < 0{,}001$, ** pour
    $p < 0{,}01$, * pour $p < 0{,}05$, et sont indicatives vu les six groupes. La valeur $p$ de
    permutation placebo suit. Le dernier bloc donne la semi-élasticité de la distribution des dividendes des ménages, sur
    base lissée des trois années antérieures au choc, estimée au niveau agrégé. Les définitions sont
    dans la table~\ref{table:definitions-variables}.\\}
    \footnotesize
    \begin{tblr}{colspec={l*{2}{X[1,c]}}}
    \toprule
     & (1) & (2) \\
     & Choc de 2013 & Choc de 2018 \\
    \midrule
    \SetCell[c=3]{l}{\textit{Sociétés privées, part de l'actif}} & & \\
    Réserves sur actif & $+0{,}0942^{*}$ & $+0{,}0332$ \\
                       & $(+2{,}2)$  & $(+0{,}6)$  \\
    Fonds propres sur actif & $+0{,}1094^{**}$ & $+0{,}0403$ \\
                            & $(+2{,}6)$ & $(+0{,}8)$ \\
     & & \\
    Placebo $p$, réserves      & $0{,}667$ & $0{,}833$ \\
    Placebo $p$, fonds propres & $0{,}333$ & $0{,}833$ \\
     & & \\
    \SetCell[c=3]{l}{\textit{Ménages, distribution agrégée}} & & \\
    Semi-élasticité des dividendes & $-0{,}40$ & $+0{,}18$ \\
     & & \\
    Effets fixes société et année & Oui & Oui \\
    Contrôle de taille            & Oui & Oui \\
    Regroupement                  & Pays & Pays \\
    \bottomrule
    \end{tblr}
\end{table}
}

\section{Robustesse}\label{sec:robustesse}

\subsection{Contrôle positif}

Le contrôle positif conforte l'interprétation. Rejoué autour du choc de 2013, notre cadre retrouve
la hausse de la rétention des bénéfices documentée par \citet{bach2023} : la rétention et les fonds
propres des sociétés françaises augmentent après la hausse d'impôt, comme le montre la colonne 2013
du tableau~\ref{table:h3-cliquet}. Cette réplication est nécessaire, car sans elle l'asymétrie ne
serait pas interprétable. Elle porte sur la rétention, non sur l'investissement : une première étude
du même choc avait rapporté une réponse de l'investissement \citep{boissel2022}, mais cette étude a
été rétractée après la mise en évidence de tendances pré-existantes \citep{bachcomment2023}, ce qui
déplace l'ancrage robuste vers l'épargne interne des sociétés, précisément la marge que nous
retrouvons.

\subsection{Contrôle négatif}

Le contrôle négatif tient. Sur les sociétés cotées, où aucun ménage résident ne pilote la
distribution, la réforme de 2018 ne fait entrer aucun capital neuf et ne relève pas
l'investissement. Un léger mouvement de la structure de financement cotée apparaît, la part des
fonds propres baissant et le levier montant, mais il provient d'un rétrécissement des fonds propres
par rachats et moindre émission, non d'un recours à la dette dont le flux ne croît pas, et il est
confondu par la baisse concomitante de l'impôt sur les sociétés et par le cycle mondial de rachats.
Il ne relève pas du canal étudié.

\subsection{Anticipation}

L'anticipation est traitée. Les distributions se reportent avant une hausse annoncée et se
retiennent avant une baisse annoncée, ce qui peut gonfler les sauts d'une année isolée. La mesure
sur base lissée, qui remplace l'année de référence par la moyenne des trois années antérieures, ne
modifie pas la conclusion : le rapport des semi-élasticités de la distribution reste d'environ
quarante-cinq pour cent, qu'on le mesure sur le flux brut d'une année sur l'autre ou sur la base
lissée. Le cliquet de distribution n'est donc pas un artefact du calendrier des versements.

\subsection{Un subside aux fonds propres chez certains contrôles}

Trois pays de contrôle subventionnent directement les fonds propres par une allocation pour fonds
propres de type \textit{ACE} : la Belgique par ses intérêts notionnels depuis 2006, l'Italie par son
aide à la croissance économique de 2011 à 2018 puis à partir de 2020, l'Espagne par sa réserve de
capitalisation depuis 2015. Un tel subside agit sur la marge même de l'hypothèse H3, la rétention et
les fonds propres, et fait de l'Espagne un pays en partie traité. Nous ajoutons donc un indicateur
d'allocation par pays et par année, fixé sur la loi et non estimé, sans retirer aucun pays. Le
cliquet de capital survit : la rétention de 2013 passe de $+0{,}0942$ à $+0{,}1061$, l'asymétrie se
renforçant d'un rapport d'environ deux fois et demie à environ trois fois et demie.\footnote{Avec
l'indicateur d'allocation, rétention 2013 $+0{,}1061$ et 2018 $+0{,}0305$, soit
$0{,}1061 / 0{,}0305 = 3{,}48$.} Comparée aux seuls contrôles exempts de tout subside, l'Allemagne et
les Pays-Bas, la rétention française de 2013 ressort à $+0{,}1337$ et le motif est encore plus
asymétrique.

\subsection{Le cycle conjoncturel ne produit pas l'asymétrie}

Le creux conjoncturel de 2013, au plus bas de la crise de la zone euro, pourrait gonfler l'écart de
cette année-là. Avec un contrôle de chômage par pays et par année, en plus de l'indicateur
d'allocation, le coefficient de 2013 reste quasi inchangé, la rétention à $+0{,}1076$. Le
resserrement apparent de l'asymétrie vient du dénominateur, c'est 2018 qui se relève légèrement.
L'argument décisif tient à la nature du contraste : le cycle ne peut pas fabriquer l'asymétrie,
puisque 2018 n'est pas une année de crise, et l'objet du papier est précisément le contraste entre
2013 et 2018, immunisé par construction contre un cycle commun aux six pays.

\subsection{Choix du pool de comparaison}

La spécification d'identification retient cinq pays comparables, économies de taille et de
structure proches, intégrées au même marché unique, car un pool valide se compose d'unités qui
auraient plausiblement suivi la même trajectoire sans la réforme. L'élargissement à des économies
dissemblables dilue le rang de la France par construction et n'informe pas sur l'identification ; il
la conforte en montrant que l'effet est propre aux économies comparables. L'inférence placebo étant
mécaniquement limitée par le petit nombre d'unités, la crédibilité repose sur la convergence de
tests indépendants, le contrôle positif, le contrôle négatif et la cohérence entre la distribution
et le capital, plutôt que sur une seule valeur $p$ de permutation.

\subsection{Interprétation du résultat nul de 2018 : un plancher déjà atteint}

Une lecture du nul de 2018 est qu'il traduit un plancher déjà atteint : la hausse de 2013 ayant
poussé le capital vers les fonds propres, la baisse de 2018 n'avait plus de capital neuf à faire
entrer. La figure~\ref{fig:capital-retenu} suit le niveau de capital retenu de la France, en écart
aux cinq pays comparables, référence 2012. La forme est celle qu'attend le mécanisme : la France n'est pas
au-dessus des pays comparables avant 2013, les deux marges sautent en 2013, puis restent élevées jusqu'en
2022 sans que la baisse de 2018 ne les ramène. Nous présentons cette lecture comme une hypothèse
corroborée et non établie. La pré-période propre ne compte que deux années avant la référence, trop
peu pour fixer la normalité d'avant 2013 ; l'inférence à six pays reste mince ; et la trajectoire est
relative aux contrôles, de sorte qu'elle établit que l'élévation est propre à la France après 2013,
non qu'elle est causée par la seule réforme de 2013 à l'exclusion de tout autre fait français de la
même année.

\begin{figure}[H]
  \centering
  \begin{tikzpicture}
    \begin{axis}[paperaxis,
      xlabel={Année}, ylabel={Écart France moins comparables (part de l'actif)},
      xtick={2010,2011,2012,2013,2014,2015,2016,2017,2018,2019,2020,2021,2022},
      xticklabel style={/pgf/number format/1000 sep=, rotate=45, anchor=east},
      legend pos=north west, legend cell align=left,
    ]
      \draw[zeroline, line width=0.6pt] (axis cs:2010,0) -- (axis cs:2022,0);
      % Deux marges du capital retenu de la France, toutes deux en bleu (France) ; le style de trait distingue.
      \addplot[franceblue, line width=1.3pt, mark=*, mark options={scale=0.7}]
        table[col sep=comma, x=year, y=retention] {figures/capital_retenu.csv};
      \addlegendentry{Réserves sur actif}
      \addplot[franceblue, line width=1.3pt, densely dashed, mark=square*, mark options={scale=0.6}]
        table[col sep=comma, x=year, y=equity] {figures/capital_retenu.csv};
      \addlegendentry{Fonds propres sur actif}
      % Les deux chocs, chacun dans sa couleur : orange la hausse de 2013, vert la baisse de 2018.
      \draw[hausse2013, dashed, line width=1pt]
        ({axis cs:2012.5,0}|-{rel axis cs:0,0}) -- ({axis cs:2012.5,0}|-{rel axis cs:0,1});
      \draw[baisse2018, dashed, line width=1pt]
        ({axis cs:2017.5,0}|-{rel axis cs:0,0}) -- ({axis cs:2017.5,0}|-{rel axis cs:0,1});
    \end{axis}
  \end{tikzpicture}
  \caption{Niveau de capital retenu de la France, référence 2012, autour des deux chocs}
  \label{fig:capital-retenu}
  \fignote{Coefficients $\tau_j$ de l'équation~\eqref{eq:event} sur le niveau de capital retenu,
  réserves sur actif et fonds propres sur actif, écart de la France aux cinq pays comparables l'année $j$,
  année de référence 2012 omise, panel étendu à 2010. Les deux séries sont en bleu, la couleur de la
  France dans tout le papier ; le trait plein est la rétention, le trait tireté les fonds propres.
  La ligne orange verticale marque la hausse d'impôt de 2013, la ligne verte la baisse de 2018. Un
  profil plat avant 2013, un saut en 2013, puis un plateau sans retour en 2018 est la forme qui
  soutient le cliquet.}
\end{figure}


\subsection{Le pic de queue ne porte pas l'identification}

On pourrait soupçonner le résultat d'absence d'entrée de capital de reposer sur le creux de la
figure~\ref{fig:entree} en 2021, année de rebond post-pandémie. Il n'en est rien. En arrêtant la
fenêtre d'estimation en 2019, avant la crise sanitaire et avant les prêts garantis par l'État, le
nul de l'hypothèse H1 tient : le coefficient de double différence sur la variation du capital
souscrit vaut $-0{,}0017$, une grandeur économiquement nulle, et l'étude d'événement est plate avant
la réforme, $+0{,}0008$ en 2014, $+0{,}0019$ en 2015 et $+0{,}0010$ en 2016, puis nulle juste après,
$-0{,}0002$ en 2018 et $-0{,}0014$ en 2019. L'absence d'entrée de capital est donc établie sur la
seule période antérieure à la pandémie et ne doit rien au mouvement des années 2020 à 2022.

\subsection{Un confond propre à la crise sanitaire : les prêts garantis par l'État}

Les années 2020 à 2022 portent une politique française étrangère au canal de la flat tax. Le
dispositif des prêts garantis par l'État, environ 143 milliards d'euros distribués à plus de 680\,000 entreprises \citep{courdescomptes2022pge},
a déversé sur les sociétés françaises une dette garantie bon marché sans équivalent de même ampleur
chez les pays comparables. La figure~\ref{fig:covid-pge} suit deux positions de bilan, la trésorerie
sur actif et la dette financière sur actif, en écart de la France aux pays comparables et référence
2017. Les deux séries sont plates jusqu'en 2019 puis bondissent conjointement en 2020, 2021 et 2022,
la trésorerie de $+0{,}031$, $+0{,}035$ et $+0{,}022$, la dette financière de $+0{,}029$, $+0{,}033$
et $+0{,}035$. Le levier total de la figure~\ref{fig:levier} ne bouge pas pour autant, car la dette
garantie nouvelle est logée en trésorerie, une réserve de liquidité, plutôt que convertie en
distribution. Ce mouvement est une corrélation descriptive autour de la pandémie, non une preuve que
la dette a financé les versements ; l'établir exigerait des données de recours au dispositif
appariées aux dividendes au niveau de la firme. Les résultats de la section~\ref{sec:resultats} n'en
dépendent pas : l'effet de la réforme est identifié sur le saut de 2018, antérieur à la pandémie, et
il tient lorsque la fenêtre s'arrête en 2019.

La poussée des distributions de 2021 ne suit d'ailleurs pas les résultats des entreprises. La
figure~\ref{fig:decouplage} place côte à côte, pour la France et en indice base 2017, les dividendes
reçus par les ménages, le profit opérationnel des sociétés non financières et l'activité. Les
dividendes retombent à 126 en 2020 puis bondissent à 187 en 2021, très au-dessus de l'excédent brut
d'exploitation des sociétés non financières, à 119, et du produit intérieur brut en volume, à 103,
l'un et l'autre tout juste revenus vers leur niveau de 2019 ; l'écart aux pays comparables va dans le
même sens, les dividendes français dépassant la moyenne des comparables de 80 points d'indice quand
le profit ne la dépasse que de 6 et l'activité d'aucun.\footnote{Indices base 2017 $= 100$, année
2021, Eurostat : dividendes reçus des ménages France $186{,}9$ et comparables $106{,}6$, soit un
écart de $186{,}9 - 106{,}6 = 80{,}3$ ; excédent brut d'exploitation des sociétés non financières
France $119{,}0$ et comparables $112{,}9$, écart $6{,}1$ ; produit intérieur brut en volume France
$102{,}6$ et comparables $103{,}0$, écart $-0{,}4$. Deux réserves de périmètre : les dividendes des
ménages incluent les versements reçus de l'étranger et des fonds, non les seuls dividendes des
sociétés non financières françaises, et l'excédent brut d'exploitation est un profit opérationnel
avant intérêts et impôts, non le résultat net distribuable.} Les dividendes reçus par les ménages se
découplent ainsi de l'activité et du profit opérationnel domestiques, un profil cohérent avec un
rattrapage après le creux de 2020 et avec le matelas de liquidité constitué par la dette garantie,
non avec un surcroît de bénéfices.

\begin{figure}[H]
  \centering
  \begin{tikzpicture}
    \begin{axis}[paperaxis,
      xlabel={Année}, ylabel={Écart France moins comparables (part de l'actif)},
      xtick={2013,2014,2015,2016,2017,2018,2019,2020,2021,2022},
      xticklabel style={/pgf/number format/1000 sep=, rotate=45, anchor=east},
      legend pos=north west, legend cell align=left,
    ]
      \draw[zeroline, line width=0.6pt] (axis cs:2013,0) -- (axis cs:2022,0);
      % Deux positions de bilan de la France, toutes deux en bleu ; le style de trait distingue.
      \addplot[franceblue, line width=1.3pt, mark=*, mark options={scale=0.7}]
        table[col sep=comma, x=year, y=coef] {figures/cash_event.csv};
      \addlegendentry{Trésorerie sur actif}
      \addplot[franceblue, line width=1.3pt, densely dashed, mark=square*, mark options={scale=0.6}]
        table[col sep=comma, x=year, y=coef] {figures/fin_debt_event.csv};
      \addlegendentry{Dette financière sur actif}
      % Reforme de 2018 en vert ; entree dans la crise sanitaire en gris.
      \draw[baisse2018, dashed, line width=1pt]
        ({axis cs:2017.5,0}|-{rel axis cs:0,0}) -- ({axis cs:2017.5,0}|-{rel axis cs:0,1});
      \draw[controlgrey, dashed, line width=1pt]
        ({axis cs:2019.5,0}|-{rel axis cs:0,0}) -- ({axis cs:2019.5,0}|-{rel axis cs:0,1});
    \end{axis}
  \end{tikzpicture}
  \caption{Trésorerie et dette financière des sociétés françaises, écart aux pays comparables, autour de la crise sanitaire}
  \label{fig:covid-pge}
  \fignote{Coefficients $\tau_j$ de l'équation~\eqref{eq:event} pour la trésorerie sur actif et pour
  la dette financière, dette à long terme plus emprunts courants, sur actif, écart de la France aux
  cinq pays comparables l'année $j$, année de référence 2017 omise. Les deux séries sont en bleu, la
  couleur de la France dans tout le papier ; le trait plein est la trésorerie, le trait tireté la
  dette financière. La ligne verte verticale marque la réforme de 2018, la ligne grise l'entrée dans
  la crise sanitaire fin 2019. Les deux séries sont plates jusqu'en 2019 puis bondissent conjointement
  en 2020, 2021 et 2022, la signature des prêts garantis par l'État. Il s'agit d'une corrélation
  descriptive, non d'une preuve que la dette a financé la distribution.}
\end{figure}


\begin{figure}[H]
  \centering
  \begin{tikzpicture}
    \begin{axis}[paperaxis,
      xlabel={Année}, ylabel={Indice, 2017 $= 100$},
      xtick={2017,2018,2019,2020,2021,2022},
      xticklabel style={/pgf/number format/1000 sep=},
      legend pos=north west, legend cell align=left,
    ]
      % Dividendes des menages : la serie focale, en bleu France.
      \addplot[franceblue, line width=1.5pt, mark=*, mark options={scale=0.8}]
        table[col sep=comma, x=year, y=dividendes] {figures/decouplage_fr.csv};
      \addlegendentry{Dividendes des ménages}
      % Profit et activite : les references, en gris ; trait plein le profit, tirete l'activite.
      \addplot[controlgrey, line width=1.1pt, mark=square*, mark options={scale=0.6}]
        table[col sep=comma, x=year, y=profit] {figures/decouplage_fr.csv};
      \addlegendentry{Profit opérationnel (EBE des SNF)}
      \addplot[controlgrey, line width=1.1pt, densely dashed, mark=triangle*, mark options={scale=0.7}]
        table[col sep=comma, x=year, y=pib] {figures/decouplage_fr.csv};
      \addlegendentry{Activité (PIB en volume)}
      % Reforme de 2018, en vert.
      \draw[baisse2018, dashed, line width=1pt]
        ({axis cs:2018,0}|-{rel axis cs:0,0}) -- ({axis cs:2018,0}|-{rel axis cs:0,1});
    \end{axis}
  \end{tikzpicture}
  \caption{Dividendes, profit et activité en France, indice base 100 en 2017}
  \label{fig:decouplage}
  \fignote{Trois séries françaises indexées à 100 en 2017 : les dividendes reçus par les ménages,
  poste D42 des comptes de secteur ; l'excédent brut d'exploitation des sociétés non financières, le
  profit opérationnel ; et le produit intérieur brut en volume, l'activité. En bleu la série focale,
  les dividendes ; en gris les deux grandeurs de référence, trait plein le profit, trait tireté
  l'activité. La ligne verte verticale marque la réforme de 2018. Les dividendes retombent en 2020
  puis bondissent en 2021 à 187, très au-dessus du profit, à 119, et de l'activité, à 103, qui ne
  font que revenir vers leur niveau de 2019. L'écart aux pays comparables en 2021 va dans le même
  sens : dividendes $+80$ points d'indice, profit $+6$, activité nul. Deux réserves de périmètre : les
  dividendes des ménages incluent les versements reçus de l'étranger et des fonds, pas seulement des
  sociétés non financières françaises, et l'excédent brut d'exploitation est un profit opérationnel
  avant intérêts et impôts, non le résultat net distribuable. La figure montre donc que les dividendes
  reçus par les ménages se découplent de l'activité et du profit opérationnel domestiques, non que les
  entreprises distribuent sans profit. Source : Eurostat.}
\end{figure}


\section{Conclusion}\label{sec:conclusion}

Ce papier demandait si une baisse de la fiscalité des dividendes fait entrer du capital dans les
entreprises ou nourrit seulement la distribution aux ménages, et si la réponse du capital à l'impôt
est symétrique. La réforme française de 2018, le prélèvement forfaitaire unique, et la hausse
d'impôt de 2013 qui l'avait précédée sur la même marge fournissent le double choc qui permet d'y
répondre.

Trois résultats ressortent, chacun répondant à une hypothèse. Sur la première, la baisse d'impôt de
2018 ne fait entrer aucun capital neuf : le coefficient sur la variation du capital souscrit,
$-0{,}0034$, est un nul estimé avec précision, et la permutation placebo place la France au plus
faible écart des six pays. Sur la deuxième, aucune substitution par la dette ni effet réel
n'apparaît ; la complémentarité comptable entre part de fonds propres et part de dettes rend les
deux coefficients égaux au signe près et tous deux nuls, de sorte qu'au point de fonds propres perdu
ne répond aucun point de dette gagné. Ces deux résultats confirment la neutralité des nouvelles vues
de l'impôt sur les dividendes, dans laquelle la fiscalité de la distribution ne déplace pas la marge
d'investissement. L'argent libéré par la réforme sort vers les ménages, dont les dividendes reçus
bondissent de 32 \% l'année même, et se loge en épargne liquide et en actions de portefeuille, non
dans le capital des sociétés privées ni dans l'immobilier. Cet écart entre la vive réponse de la
distribution et l'absence de réponse réelle rejoint ce que \citet{chetty2005} et \citet{yagan2015}
observent sur la baisse américaine de 2003.

Le résultat central est le troisième. En plaçant la hausse de 2013 et la baisse de 2018 dans un
cadre d'estimation unique, nous mettons au jour un cliquet : sur la distribution, la baisse ne défait
que 45 \% de l'effet de la hausse ; sur l'allocation du capital des sociétés privées, elle n'en
ramène qu'environ un tiers, la réallocation revenant dans un rapport de un à près de trois sur les
réserves comme sur les fonds propres. Le capital part plus facilement qu'il ne revient. La
contribution du papier tient à ce contraste placé dans un cadre unique, qui fait de l'asymétrie
elle-même l'objet de mesure là où la littérature n'observait qu'un choc isolé, et à sa cohérence sur
deux niveaux indépendants, la distribution agrégée des ménages et l'allocation du capital des
sociétés.

L'implication pour l'évaluation de la fiscalité du capital est directe et forte. Une baisse d'impôt
symétrique à une hausse passée ne défait pas ses effets : les déformations qu'une hausse imprime à
l'allocation du capital persistent après la baisse qui lui succède, parce que les coûts d'ajustement
et les motifs de rétention rendent la sortie des fonds plus réactive que leur retour. Un calcul
d'évaluation qui suppose la réversibilité surestime donc le gain d'une baisse et sous-estime le coût
durable d'une hausse.

La portée du résultat est bornée par le cadre d'identification. L'inférence repose sur six pays, ce
qui limite la puissance de la permutation placebo, et la trajectoire du capital retenu est mesurée
relativement aux contrôles, de sorte qu'elle établit que l'élévation postérieure à 2013 est propre à
la France sans l'imputer à la seule réforme de 2013 à l'exclusion de tout autre fait français de la
même année. Ces limites invitent à répliquer le contraste des deux chocs sur d'autres pays ayant
parcouru la même marge fiscale dans les deux sens, que les données microéconomiques de sociétés
privées rendent désormais accessibles.


\clearpage
\bibliography{references}

\clearpage
\appendix
\section*{Annexe}
\addcontentsline{toc}{section}{Annexe}

{
\small
\begin{longtblr}[
  caption={Description et définitions des variables},
  label={table:definitions-variables}
]{
  width=\linewidth,
  colspec={X[2,l]X[6,j]X[2,l]},
  column{3}={font=\itshape},
  row{1}={font=\bfseries},
}
\toprule
Variable & Définition & Source \\
\midrule \addlinespace[6pt]
\SetCell[c=3]{l}{\textbf{\textit{Variables expliquées}}} & & \\
\midrule
Variation du capital $g_{it}$ & Variation annuelle du logarithme du capital souscrit et libéré,
  $g_{it} = \ln(1 + K_{it}) - \ln(1 + K_{i,t-1})$, où $K_{it}$ est le capital souscrit. & Orbis, Bureau van Dijk \\
Hausse du capital $a_{it}$ & Indicatrice valant un si le capital souscrit augmente d'une année sur
  l'autre, zéro sinon. & Orbis, Bureau van Dijk \\
Ratio de fonds propres $e_{it}$ & Capitaux propres sur total de l'actif, $e_{it} = E_{it} / A_{it}$. & Orbis, Bureau van Dijk \\
Levier $\ell_{it}$ & Total des dettes sur total de l'actif. & Orbis, Bureau van Dijk \\
Croissance de la dette & Variation du logarithme de la dette financière, long terme plus emprunts
  courants. & Orbis, Bureau van Dijk \\
Investissement & Variation du logarithme des immobilisations corporelles. & Orbis, Bureau van Dijk \\
Productivité du travail & Logarithme du chiffre d'affaires par salarié. & Orbis, Bureau van Dijk \\
Réserves sur actif $r_{it}$ & Capitaux propres moins capital souscrit, sur total de l'actif,
  $r_{it} = (E_{it} - K_{it}) / A_{it}$. & Orbis, Bureau van Dijk \\
Dividendes des ménages $Y_{ct}$ & Logarithme des dividendes distribués reçus par les ménages du pays
  $c$ l'année $t$, poste des comptes de secteur. & Eurostat, Office statistique de l'Union européenne \\
\midrule \addlinespace[6pt]
\SetCell[c=3]{l}{\textbf{\textit{Traitement et contrôle}}} & & \\
\midrule
France $\times$ après réforme & Produit de l'indicatrice France et de l'indicatrice d'année
  postérieure au choc ; le régresseur de double différence de l'équation~\eqref{eq:did}. & Construction \\
Taille $s_{it}$ & Logarithme du total de l'actif, $s_{it} = \ln(1 + A_{it})$. & Orbis, Bureau van Dijk \\
\bottomrule
\end{longtblr}
}

\bigskip

{
\footnotesize
\begin{longtblr}[
  caption={Statistiques descriptives},
  label={table:stats-descriptives}
]{
  width=\linewidth,
  colspec={l*{6}{X[1,c]}},
  row{1}={font=\bfseries},
}
\toprule
Variable & Obs & Moyenne & Écart-type & P25 & Médiane & P75 \\
\midrule \addlinespace[6pt]
\SetCell[c=7]{l}{\textbf{\textit{Variables expliquées}}} & & & & & & \\
\midrule
Variation du capital $g_{it}$   & 14\,039\,102 & $0{,}01$  & $0{,}08$ & $0{,}00$   & $0{,}00$  & $0{,}00$ \\
Hausse du capital $a_{it}$      & 14\,039\,102 & $0{,}02$  & $0{,}15$ & $0{,}00$   & $0{,}00$  & $0{,}00$ \\
Ratio de fonds propres $e_{it}$ & 16\,150\,854 & $0{,}06$  & $2{,}01$ & $0{,}09$   & $0{,}36$  & $0{,}68$ \\
Levier $\ell_{it}$              & 16\,154\,361 & $0{,}94$  & $2{,}01$ & $0{,}32$   & $0{,}64$  & $0{,}91$ \\
Croissance de la dette          & 14\,039\,102 & $-0{,}08$ & $3{,}14$ & $-0{,}13$  & $0{,}00$  & $0{,}00$ \\
Investissement                  & 13\,165\,380 & $-0{,}03$ & $1{,}52$ & $-0{,}24$  & $-0{,}03$ & $0{,}00$ \\
Productivité du travail         & 5\,445\,689  & $11{,}30$ & $0{,}91$ & $10{,}76$  & $11{,}29$ & $11{,}85$ \\
Réserves sur actif $r_{it}$     & 16\,150\,854 & $-0{,}27$ & $3{,}00$ & $-0{,}01$  & $0{,}19$  & $0{,}49$ \\
\midrule \addlinespace[6pt]
\SetCell[c=7]{l}{\textbf{\textit{Contrôle}}} & & & & & & \\
\midrule
Taille $s_{it}$                 & 16\,169\,607 & $12{,}28$ & $1{,}57$ & $11{,}45$  & $12{,}47$ & $13{,}34$ \\
\bottomrule
\end{longtblr}
}

\medskip
\noindent{\footnotesize\textit{Notes.} Panel de sociétés privées de la France et de cinq pays comparables,
Orbis, Bureau van Dijk, comptes annuels de 2013 à 2022. Les variables continues sont winsorisées aux
premier et quatre-vingt-dix-neuvième centiles. Le minimum et le maximum ne sont pas reportés ; la
dispersion se lit sur l'écart-type et l'intervalle interquartile. Les effectifs varient selon la
variable, les variations annuelles et la productivité étant définies sur les seules années-sociétés
qui portent les lignes de compte requises.}

\bigskip

{
\begin{table}[ht]
    \centering
    \caption{Chocs fiscaux exploités et stratégie d'identification}\label{table:identification}
    \tablenotes{Variation identifiante du papier, détaillée dans la section~\ref{sec:modele}. Pays
    traité : la France, deux chocs de sens opposé sur l'imposition des dividendes. Pays de contrôle :
    régime forfaitaire du capital stable, année d'instauration entre parenthèses, servant de
    contre-factuel sur toute la fenêtre 2010 à 2022. Le tiret marque une case sans objet. Les dates
    des textes de loi sont à confirmer avant diffusion. Sources : loi de finances pour 2018
    \citep{loidefinances2018}, base fiscale de l'OCDE, \citet{sorensen1994} pour la lignée de
    l'impôt dual.\\}
    \footnotesize
    \begin{tblr}{colspec={X[1.2,l] X[3,l] X[0.9,c] X[0.7,c] X[1,c] X[1,c]}}
    \toprule
    Pays & Réforme et instrument légal & Sens & Année & Fenêtre avant & Fenêtre après \\
    \midrule
    \SetCell[c=6]{l}{\textit{Pays traité}} & & & & & \\
    France & Loi de finances 2013, retour des dividendes au barème progressif & Hausse & 2013 & 2010--2012 & 2013--2016 \\
    France & Loi de finances 2018, article 28, prélèvement forfaitaire unique à 30 \% & Baisse & 2018 & 2014--2017 & 2018--2022 \\
    \midrule
    \SetCell[c=6]{l}{\textit{Pays de contrôle, régime forfaitaire du capital stable (année d'instauration)}} & & & & & \\
    Allemagne & \textit{Abgeltungsteuer}, prélèvement forfaitaire à 26,375 \% (2009) & \SetCell[c=4]{l}{Contre-factuel sur 2010--2022} & & & \\
    Pays-Bas  & \textit{Box 3}, imposition forfaitaire du patrimoine (2001) & \SetCell[c=4]{l}{Contre-factuel sur 2010--2022} & & & \\
    Espagne   & Base de l'épargne, barème dédié au capital (2007) & \SetCell[c=4]{l}{Contre-factuel sur 2010--2022} & & & \\
    Belgique  & Précompte mobilier forfaitaire sur dividendes & \SetCell[c=4]{l}{Contre-factuel sur 2010--2022} & & & \\
    Italie    & \textit{Cedolare secca}, prélèvement forfaitaire sur dividendes & \SetCell[c=4]{l}{Contre-factuel sur 2010--2022} & & & \\
    \bottomrule
    \end{tblr}
\end{table}
}


\end{document}
